\documentclass{article}
\usepackage{lmodern}
\usepackage{amsmath}
\usepackage{listings}
\usepackage{fullpage}
\usepackage{parskip}
\usepackage{xcolor}
\lstset{
	basicstyle=\ttfamily,
	columns=fullflexible,
	frame=single,
	breaklines=true,
	postbreak=\mbox{\textcolor{red}{$\hookrightarrow$}\space},
}
\begin{document}
\title{Experiment `p\_6\_2\_h\_is\_sorted\_array' Results}
\author{\today}
\date{}
\maketitle
\textbf{Experiment outcome: }SUCCESS
\\\textbf{Bad responses: }0
\\\textbf{Responses containing}~\texttt{assume}~\textbf{: }0
\\\textbf{Responses containing}~\texttt{expect}~\textbf{: }0
\\\textbf{Resolution attempts: }1
\\\textbf{Hard fails (resolution): }0
\\\textbf{Soft fails (resolution): }0
\\\textbf{Verification attempts: }1
\section*{Problem Specification}
\textbf{Problem name: }p\_6\_2\_h\_is\_sorted\_array
\\\textbf{Natural language statement: }null
\\\textbf{Method signature: }p\_6\_2\_h\_is\_sorted\_array(inputs: array<int>) returns (is\_sorted: bool)
\subsection*{Ensures}
\begin{itemize}
\item \texttt{is\_sorted == (forall i :: 0 <= i < inputs.Length - 1 ==> inputs[i] <= inputs[i + 1])}
\end{itemize}
\clearpage
\section*{GenAI interactions}
Below you will find all interactions between the `user' (program) and the `assistant' (OpenAI).
\subsection*{Program $\rightarrow$ GenAI}
\begin{lstlisting}
You are given the following task to perform in Dafny:

Write a method to satisfy the below requirements.

The signature should be:

p_6_2_h_is_sorted_array(inputs: array<int>) returns (is_sorted: bool)

The method should respect the following contract:

ensures is_sorted == (forall i :: 0 <= i < inputs.Length - 1 ==> inputs[i] <= inputs[i + 1])

Produce and show only the Dafny body of this method, including the curly braces that surround it. Do not show the signature nor contract. You must not use 'assume' nor 'expect' anywhere in your code.
\end{lstlisting}
\subsection*{GenAI $\rightarrow$ Program}
\textbf{System fingerprint: }null
\\\textbf{ID: }chatcmpl-DnO4t0Vwge1bCJTI3vgBJJ52ZtgaA
\\\textbf{Model: }gpt-5.5-2026-04-23
\\\textbf{Created at: }1780662999
\\\textbf{Finish reason: }stop
\\\textbf{Usage: }296 tokens in, and 3530 tokens out
\begin{lstlisting}
{
  var i := 0;
  while i < inputs.Length - 1
    invariant 0 <= i <= inputs.Length
    invariant forall j :: 0 <= j < i && j + 1 < inputs.Length ==> inputs[j] <= inputs[j + 1]
  {
    if inputs[i] > inputs[i + 1] {
      is_sorted := false;
      assert 0 <= i < inputs.Length - 1;
      assert !(forall j :: 0 <= j < inputs.Length - 1 ==> inputs[j] <= inputs[j + 1]) by {
        if forall j :: 0 <= j < inputs.Length - 1 ==> inputs[j] <= inputs[j + 1] {
          assert inputs[i] <= inputs[i + 1];
          assert false;
        }
      }
      return;
    }

    var old_i := i;
    assert inputs[old_i] <= inputs[old_i + 1];
    i := i + 1;
    assert forall j :: 0 <= j < i && j + 1 < inputs.Length ==> inputs[j] <= inputs[j + 1] by {
      forall j | 0 <= j < i && j + 1 < inputs.Length
        ensures inputs[j] <= inputs[j + 1]
      {
        if j == old_i {
          assert inputs[j] <= inputs[j + 1];
        } else {
          assert j < old_i;
          assert inputs[j] <= inputs[j + 1];
        }
      }
    }
  }

  is_sorted := true;
  assert forall j :: 0 <= j < inputs.Length - 1 ==> inputs[j] <= inputs[j + 1] by {
    forall j | 0 <= j < inputs.Length - 1
      ensures inputs[j] <= inputs[j + 1]
    {
      assert !(i < inputs.Length - 1);
      assert inputs.Length - 1 <= i;
      assert j < i;
      assert j + 1 < inputs.Length;
      assert inputs[j] <= inputs[j + 1];
    }
  }
}
\end{lstlisting}
\section*{Final Program}
\begin{lstlisting}
method p_6_2_h_is_sorted_array(inputs: array<int>) returns (is_sorted: bool)
	ensures is_sorted == (forall i :: 0 <= i < inputs.Length - 1 ==> inputs[i] <= inputs[i + 1])
{
  var i := 0;
  while i < inputs.Length - 1
    invariant 0 <= i <= inputs.Length
    invariant forall j :: 0 <= j < i && j + 1 < inputs.Length ==> inputs[j] <= inputs[j + 1]
  {
    if inputs[i] > inputs[i + 1] {
      is_sorted := false;
      assert 0 <= i < inputs.Length - 1;
      assert !(forall j :: 0 <= j < inputs.Length - 1 ==> inputs[j] <= inputs[j + 1]) by {
        if forall j :: 0 <= j < inputs.Length - 1 ==> inputs[j] <= inputs[j + 1] {
          assert inputs[i] <= inputs[i + 1];
          assert false;
        }
      }
      return;
    }

    var old_i := i;
    assert inputs[old_i] <= inputs[old_i + 1];
    i := i + 1;
    assert forall j :: 0 <= j < i && j + 1 < inputs.Length ==> inputs[j] <= inputs[j + 1] by {
      forall j | 0 <= j < i && j + 1 < inputs.Length
        ensures inputs[j] <= inputs[j + 1]
      {
        if j == old_i {
          assert inputs[j] <= inputs[j + 1];
        } else {
          assert j < old_i;
          assert inputs[j] <= inputs[j + 1];
        }
      }
    }
  }

  is_sorted := true;
  assert forall j :: 0 <= j < inputs.Length - 1 ==> inputs[j] <= inputs[j + 1] by {
    forall j | 0 <= j < inputs.Length - 1
      ensures inputs[j] <= inputs[j + 1]
    {
      assert !(i < inputs.Length - 1);
      assert inputs.Length - 1 <= i;
      assert j < i;
      assert j + 1 < inputs.Length;
      assert inputs[j] <= inputs[j + 1];
    }
  }
}
\end{lstlisting}
\section*{Total Token Usage}
\textbf{Input tokens: }296
\\\textbf{Output tokens: }3530
\\\textbf{Reasoning tokens: }2997
\\\textbf{Sum of `total tokens': }3826
\section*{Experiment Timings}
\textbf{Overall Experiment} started at 1780662999129, ended at 1780663093099, lasting 93970ms (93.97 seconds)
\\\textbf{Iteration \#1} started at 1780662999129, ended at 1780663093099, lasting 93970ms (93.97 seconds)
\end{document}
